%%%%%%%%%%%%%%%%%%%%
%%% Mise en page %%%
%%%%%%%%%%%%%%%%%%%%

\newcommand\titre[1]{\begin{center} \textbf{\LARGE #1} \end{center}}

\newcommand*\circled[1]{\tikz[baseline=(char.base)]{
            \node[shape=circle,draw,inner sep=2pt] (char) {#1};}}

            
%%%%%%%%%%%%%%%%%%
%%%% Ensembles %%%
%%%%%%%%%%%%%%%%%%

\newcommand\CC{%
        {\mathbb C}} % ----- Ensemble C

\newcommand\RR{%
        {\mathbb R}} % ----- Ensemble R

\newcommand\QQ{%
        {\mathbb Q}} % ----- Ensemble Q

\newcommand\ZZ{%
        {\mathbb Z}} % ----- Ensemble Z

\newcommand\NN{%
        {\mathbb N}} % ----- Ensemble N

\newcommand\PP{%
        {\mathbb P}} % ----- P probas

\newcommand\EE{%
        {\mathbb E}} % ----- E espérance

\newcommand\VV{%
        {\mathbb V}} % ----- V variance


\newcommand\nN{%
        {n\in\mathbb N}} % ----- n dans N

\newcommand\xR{%
        {x\in\mathbb R}} % ----- x dans R

\newcommand\zC{%
        {z\in\mathbb C}} % ----- z dans C


%%%%%%%%%%%%%%%%
%%% Symboles %%%
%%%%%%%%%%%%%%%%

\newcommand\dur{%
        {(\star)}} % ----- Exo difficile



\newcommand\limx{%
        {\lim\limits_{x\to +\infty}}\,} % ----- Limite x en plus infini

\newcommand\limn{%
        {\lim\limits_{n\to +\infty}}\,} % ----- Limite n en plus infini

\newcommand\limite[1]{%
        {\lim\limits_{\substack{#1}}}\,} % ----- Limite X verx Y

\newcommand\Ccal{%
        {\mathscr{C}}} % ----- C calligraphique

\newcommand\Dcal{%
        {\mathscr{D}}} % ----- D calligraphique

\newcommand\Pcal{%
        {\mathscr{P}}} % ----- P calligraphique

\newcommand\Rcal{%
        {\mathscr{R}}} % ----- R calligraphique

\newcommand\Fcal{%
        {\mathscr{F}}} % ----- F calligraphique

\newcommand\Scal{%
        {\mathscr{S}}} % ----- S calligraphique

\newcommand\Acal{%
        {\mathscr{A}}} % ----- A calligraphique

\newcommand\Vcal{%
        {\mathscr{V}}} % ----- V calligraphique



%%%%%%%%%%%%%%%%%%%%%%%%%
%%% Accolade à gauche %%%
%%%%%%%%%%%%%%%%%%%%%%%%%

\newcommand{\accol}[1]{%
    {\left\{
    \renewcommand{\arraystretch}{0.8}
    \begin{array}{l}
        #1
    \end{array}
    \right. }}


\newcommand{\accolsized}[2]{%
    {\left\{
    \renewcommand{\arraystretch}{#2}
    \begin{array}{l}
        #1
    \end{array}
    \right. }}



%%%%%%%%%%%%%%%%%%%%%%
%%% Calculs usuels %%%
%%%%%%%%%%%%%%%%%%%%%%

\newcommand\somme[2]{%
        {\displaystyle\sum_{#1}^{#2}}}    % ----- Somme

\newcommand\sommek[2]{%
        {\displaystyle\sum_{k=#1}^{#2}}}    % ----- Somme k

\newcommand\inte[2]{%
        {\displaystyle\int_{#1}^{#2}}}   % ----- Intégrale

\newcommand\dd{%
        {\,\,\mathrm{d}}}       % ----- d différentiel

\newcommand\floor[1]{%
        {\lfloor #1 \rfloor}}   % ----- Floor

\newcommand\ceil[1]{%
        {\lceil #1 \rceil}}     % ----- Ceiling

\newcommand\abs[1]{%
        {\left\lvert #1 \right\rvert}}     % ----- Absolute value
        
\newcommand\pgcd{%
        {\mathrm{pgcd}}}       % ----- pgcd

\newcommand\ppcm{%
        {\mathrm{ppcm}}}       % ----- ppcm




%%%%%%%%%%%%%%
%%% Suites %%%
%%%%%%%%%%%%%%

\newcommand{\defsuite}[2]{%
    \renewcommand{\arraystretch}{0.8}
    \begin{cases}
        #1 \\
        u_{n+1}= #2 \\
    \end{cases}
    }                   % ----- Définition d'une suite (Un) par récurrence

\newcommand\Un{%
        {(u_{n})_{n\in\mathbb N}}} % ----- (Un)_n in N



%%%%%%%%%%%%%%%%%
%%% Géométrie %%%
%%%%%%%%%%%%%%%%%

\newcommand\vvec[1]{\vbox{\halign{##\cr\tiny\textbf{\rightarrowfill}\cr\noalign{\nointerlineskip\vskip2pt}$#1\mskip2mu$\cr}}} % ----- Vecteur

\newcommand\vcoord[2]{\renewcommand{\arraystretch}{0.6}\vvec{#1}\hspace{-2pt}\matrice{#2}} % ----- Coordonnées vecteur

\newcommand\Oij{\left(O\,;\,\vvec{i}\,;\,\vvec{j}\right)} % ----- (O;i;j)

\def\Oijk{\left(O\,;\,\Vec{i}\,;\,\Vec{j}\,;\,\Vec{k}\right)} % ---- (O;i;j;k)

\newcommand\norme[1]{%
        {\left\lVert #1 \right\rVert}}     % ----- Norme

\newcommand\normev[1]{%
        {\left\lVert \vvec{#1} \right\rVert}}     % ----- Norme vecteur

\newcommand\scal[2]{%
    {\vvec{#1}\cdot\vvec{#2}}} % ----- Produit scalaire



%%%%%%%%%%%%%%%%%%%%
%%% Probabilités %%%
%%%%%%%%%%%%%%%%%%%%

\newcommand\barre[1]{%
    {\overline{#1}}} % ----- A barre



%%%%%%%%%%%%%%%%%%%%%%%%%
%%% Nombres complexes %%%
%%%%%%%%%%%%%%%%%%%%%%%%%

\newcommand\Ouv{(O\,;\,\vvec{u}\,;\,\vvec{v})} % ----- (O;u;v)


\newcommand\Reel{%
    {\mathfrak{Re}}} % ----- Partie réelle

\newcommand\Imag{%
    {\mathfrak{Im}}} % ----- Partie imaginaire

\newcommand{\defsuitez}[2]{%
    {\left \{
    \renewcommand{\arraystretch}{0.7}
    \begin{array}{l}
        #1 \\
        z_{n+1}= #2 \\
    \end{array}
    \right. }} % ----- Définition d'une suite complexe (Zn)n par récurrence



%%%%%%%%%%%%%%%%
%%% Matrices %%%
%%%%%%%%%%%%%%%%

\newcommand\matrice[1]{{%
    \renewcommand{\arraystretch}{0.7}
    \arraycolsep=0.7\arraycolsep\ensuremath{\begin{pmatrix}#1\end{pmatrix}}}}

